\documentclass[12pt]{article}

% Core math
\usepackage{amsmath,amssymb,amsthm}
\usepackage{mathtools}

% Diagrams
\usepackage{tikz-cd}
\usepackage{tikz}

% References
\usepackage{hyperref}
\usepackage[nameinlink,noabbrev]{cleveref}

% Graphics and tables
\usepackage{graphicx}
\usepackage{booktabs}
\usepackage{longtable}
\usepackage{array}

% Page layout
\usepackage[margin=1in]{geometry}
\usepackage{microtype}
\usepackage{enumitem}

% GrokRxiv sidebar
\usepackage{everypage}
\usepackage{xcolor}

\hypersetup{
  colorlinks=true,
  linkcolor=blue!55!black,
  citecolor=green!40!black,
  urlcolor=blue!60!black,
  pdftitle={Uniform Locality for Condensed Families of Quantum Lattice Interactions},
  pdfauthor={Matthew Long and the YonedaAI Collaboration}
}

\newtheorem{theorem}{Theorem}[section]
\newtheorem{proposition}[theorem]{Proposition}
\newtheorem{lemma}[theorem]{Lemma}
\newtheorem{corollary}[theorem]{Corollary}
\theoremstyle{definition}
\newtheorem{definition}[theorem]{Definition}
\newtheorem{example}[theorem]{Example}
\newtheorem{counterexample}[theorem]{Counterexample}
\theoremstyle{remark}
\newtheorem{remark}[theorem]{Remark}

\newcommand{\Alg}{\mathcal A}
\newcommand{\Int}{\mathcal I}
\newcommand{\Fin}{\mathcal P_0}
\newcommand{\supp}{\operatorname{supp}}
\newcommand{\diam}{\operatorname{diam}}
\newcommand{\dist}{\operatorname{dist}}
\newcommand{\Cont}{\operatorname{Cont}}
\newcommand{\Ham}{\mathfrak{Ham}}
\newcommand{\Gap}{\mathfrak{Gap}}
\newcommand{\id}{\mathrm{id}}
\newcommand{\EST}{\textsc{Established}}
\newcommand{\PROP}{\textsc{Proposed}}
\newcommand{\OBS}{\textsc{Obstructed}}
\newcommand{\OPEN}{\textsc{Open}}

\definecolor{grokgray}{RGB}{110,110,110}
\AddEverypageHook{%
  \ifnum\value{page}=1
    \begin{tikzpicture}[remember picture,overlay]
      \node[
        rotate=90,
        anchor=south,
        font=\Large\sffamily\bfseries\color{grokgray},
        inner sep=0pt
      ] at ([xshift=38pt,yshift=0.52\paperheight]current page.south west)
      {GrokRxiv:2026.07.locality-condensed-families\quad
       [\,math-ph\,]\quad
       14 Jul 2026};
    \end{tikzpicture}
  \fi
}

\title{Uniform Locality for Condensed Families of Quantum Lattice Interactions}
\author{Matthew Long \\
\textit{The YonedaAI Collaboration} \\
\textit{YonedaAI Research Collective} \\
Chicago, IL \\
\texttt{matthew@yonedaai.com} $\cdot$ \url{https://yonedaai.com}}
\date{14 July 2026}

\begin{document}

\maketitle

\begin{abstract}
We isolate the analytic input needed to place parameterized quantum lattice
interactions in a condensed mathematical setting.  The decisive condition is
not profiniteness of the parameter object.  It is a bound on an interaction
$F$-norm that is uniform in the parameter.  Under the usual summability and
convolution conditions on $F$, such a bound gives finite-volume
Lieb-Robinson estimates with constants independent of both the parameter and
the volume.  It also gives thermodynamic dynamics, uniformly on compact time
intervals.  A separate local $F$-continuity condition yields point-norm
continuity of the infinite-volume dynamics on local observables.  We package
these hypotheses into a functor on profinite sets and prove its elementary
finite-cover descent property.  This packaging is categorical.  It supplies
none of the propagation estimates by itself.

Two counterexamples mark the boundary of the result.  First, compact
parameterization and pointwise locality do not imply a uniform interaction
bound.  Second, even local parameter continuity together with a uniform
locality bound and a positive gap in every fiber does not imply a gap uniform
in the family.  The latter failure is exhibited by an on-site defect whose
strength tends to zero while its position escapes to infinity.  We state all
results with exact hypotheses, distinguish established analytic theorems from
the proposed condensed organization, and include finite executable checks in
Haskell.  The outcome is a precise first arrow in the program from local
quantum interactions to a condensed moduli stack of Hamiltonians.
\end{abstract}

\tableofcontents

\section{Introduction}

Topological phases are usually organized by paths of gapped Hamiltonians,
stable equivalence, operator algebras, or effective field theories.  Each of
these viewpoints begins with a more elementary question.  Which families of
lattice interactions define controlled dynamics in the first place?  A
classification of phases cannot be more reliable than its locality
hypotheses.

Condensed mathematics suggests testing a moduli object on profinite sets.  If
$S$ is profinite, an $S$-point should describe an $S$-parameterized family of
interactions.  This is attractive for disorder spaces such as
$D^{\mathbb Z^d}$ with $D$ finite.  It is also potentially misleading.  The
sheaf condition glues compatible data over a finite cover.  It does not bound
an infinite sum of interaction terms, prove a convolution estimate, or
construct thermodynamic time evolution.

The purpose of this paper is to separate these jobs.  The analytic layer is
provided by interaction Banach spaces and Lieb-Robinson theory.  The
parameter layer records local continuity and a common norm bound.  The
condensed layer records functoriality in a profinite test object and descent
for finite jointly surjective covers.

Our main theorem is a uniform-family corollary of the standard
Lieb-Robinson and thermodynamic-limit theorems.  Let $(\Gamma,d)$ be a
countable metric space, let $F$ satisfy uniform summability and a convolution
condition, and let $(\Phi_s)_{s\in S}$ be a family of interactions such that

\begin{equation}
  \sup_{s\in S}\|\Phi_s\|_F\le M<\infty.
  \label{eq:intro-uniform}
\end{equation}

Then the propagation constants depend on $F$ and $M$, not on $s$ or the
finite volume.  The corresponding thermodynamic dynamics exist with uniform
finite-volume convergence on compact time intervals.  If the family is
locally continuous in the sense defined below, then
$s\mapsto\tau_t^s(A)$ is norm-continuous for each local observable $A$.

The compactness of $S$ is useful only after a topology has been specified.
For example, continuity of $S\to\mathcal B_F$ implies
\cref{eq:intro-uniform} when $S$ is compact.  Product-topology disorder
families are generally not continuous in this global Banach norm.  They may
still satisfy the common bound directly and may be continuous on every local
restriction.  Keeping these two requirements separate is essential.

The paper uses four claim labels.

\begin{itemize}[leftmargin=2.3em]
\item \EST{} marks a cited theorem or a direct corollary with proof.
\item \PROP{} marks the categorical organization introduced here.
\item \OBS{} marks a statement ruled out by a counterexample.
\item \OPEN{} marks a comparison or extension not proved here.
\end{itemize}

The central claim of this paper is \EST{}.  The condensed moduli functor is
\PROP{}.  The assertion that profiniteness by itself produces uniform
locality is \OBS{}.

\subsection{Main contributions}

The contributions are deliberately narrow.

\begin{enumerate}[leftmargin=2.3em]
\item We fix a single interaction norm and state the corresponding
Lieb-Robinson estimate with its constants.
\item We prove a uniform-family theorem for finite-volume propagation and
thermodynamic dynamics.
\item We formulate a local $F$-topology suitable for product-type disorder
families and prove parameter continuity of dynamics under a common tail
bound.
\item We define a profinite functor of uniformly local interactions and prove
finite-cover descent.
\item We give counterexamples separating pointwise locality, uniform
locality, local continuity, and uniform gappedness.
\item We provide a finite Haskell companion that computes $F$-function
diagnostics, Lieb-Robinson envelopes, and finite witnesses for the logical
difference between pointwise and uniform conditions.
\end{enumerate}

We do not prove a spectral gap, construct a phase spectrum, or identify a
microscopic lattice category with an effective field theory.  Those problems
begin after the analytic foundation established here.

\subsection{Relation to the literature}

The norm and propagation estimates follow the framework developed by
Nachtergaele and Sims and its detailed treatment by Nachtergaele, Sims, and
Young \cite{NachtergaeleSims2010,NSY2019}.  The thermodynamic-limit argument
uses a uniform Cauchy estimate on local observables.  Parameter dependence is
an application of continuity and local $F$-convergence estimates in
\cite{NSY2019}.  Spectral flow is discussed only to identify which output of
this paper later phase-classification work needs \cite{BMNS2012}.  No theorem
from condensed mathematics is used to prove an analytic estimate.

\section{Quantum lattice systems and interaction norms}

\subsection{The quasi-local algebra}

\begin{definition}[Quantum lattice system]
Let $(\Gamma,d)$ be a countable metric space.  To every $x\in\Gamma$ assign a
finite-dimensional complex Hilbert space $\mathcal H_x$.  For a finite set
$\Lambda\Subset\Gamma$, define

\begin{equation}
  \mathcal H_\Lambda=\bigotimes_{x\in\Lambda}\mathcal H_x,
  \qquad
  \Alg_\Lambda=\mathcal B(\mathcal H_\Lambda).
\end{equation}

If $\Lambda_0\subseteq\Lambda$, identify $A\in\Alg_{\Lambda_0}$ with
$A\otimes 1_{\Lambda\setminus\Lambda_0}$.  The local algebra and quasi-local
algebra are

\begin{equation}
  \Alg^{\mathrm{loc}}_\Gamma
  =\bigcup_{\Lambda\Subset\Gamma}\Alg_\Lambda,
  \qquad
  \Alg_\Gamma
  =\overline{\Alg^{\mathrm{loc}}_\Gamma}^{\|\cdot\|}.
\end{equation}
\end{definition}

Finite-dimensional on-site spaces keep the presentation focused.  The
Lieb-Robinson estimate permits certain unbounded on-site Hamiltonians after
passing to an interaction picture, but unbounded interaction terms require
additional hypotheses and are outside our scope.

\begin{definition}[Interaction]
An interaction is a map

\begin{equation}
  \Phi:\Fin(\Gamma)\longrightarrow\Alg^{\mathrm{loc}}_\Gamma
\end{equation}

such that $\Phi(Z)=\Phi(Z)^*\in\Alg_Z$ for every finite $Z$.  Its
finite-volume Hamiltonian is

\begin{equation}
  H^\Phi_\Lambda=\sum_{Z\subseteq\Lambda}\Phi(Z).
\end{equation}
\end{definition}

The associated finite-volume Heisenberg dynamics is

\begin{equation}
  \tau^{\Lambda,\Phi}_t(A)
  =e^{itH^\Phi_\Lambda}Ae^{-itH^\Phi_\Lambda}.
\end{equation}

All finite-volume expressions are well defined.  The substantive issue is to
control them independently of $\Lambda$.

\subsection{F-functions}

\begin{definition}[F-function]
A non-increasing function $F:[0,\infty)\to(0,\infty)$ is an $F$-function on
$(\Gamma,d)$ if

\begin{equation}
  \|F\|
  :=\sup_{x\in\Gamma}\sum_{y\in\Gamma}F(d(x,y))<\infty
  \label{eq:F-sum}
\end{equation}

and

\begin{equation}
  C_F
  :=\sup_{x,y\in\Gamma}
  \frac{1}{F(d(x,y))}
  \sum_{z\in\Gamma}F(d(x,z))F(d(z,y))<\infty.
  \label{eq:F-conv}
\end{equation}
\end{definition}

The first condition says that tails are uniformly summable.  The second says
that one convolution does not degrade the spatial profile beyond a constant.
Repeated commutator estimates then produce powers of $C_F$ without changing
the form of the decay function.

\begin{example}[Polynomial and exponential weights]
On $\mathbb Z^\nu$ with the $\ell^1$ metric, for every $\epsilon>0$,

\begin{equation}
  F(r)=\frac{1}{(1+r)^{\nu+\epsilon}}
\end{equation}

is an $F$-function.  If $g$ is non-decreasing and subadditive, then
$F_g(r)=e^{-g(r)}F(r)$ is again an $F$-function.  In particular,

\begin{equation}
  F_a(r)=e^{-ar}F(r),\qquad a>0,
\end{equation}

produces the familiar exponential light-cone form.
\end{example}

\begin{definition}[Interaction F-norm]
For an interaction $\Phi$, set

\begin{equation}
  \|\Phi\|_F
  =\sup_{x,y\in\Gamma}
  \frac{1}{F(d(x,y))}
  \sum_{\substack{Z\Subset\Gamma\\x,y\in Z}}
  \|\Phi(Z)\|.
  \label{eq:F-norm}
\end{equation}

Let $\mathcal B_F$ be the vector space of interactions with finite
$F$-norm.
\end{definition}

The $x=y$ part controls the total interaction incident on one site.  The
$x\ne y$ part controls simultaneous incidence at separated sites.  This is
stronger than asking each individual term to decay with its diameter.

\begin{lemma}[Pair bound]
If $\Phi\in\mathcal B_F$, then for all $x,y\in\Gamma$,

\begin{equation}
  \sum_{Z\ni x,y}\|\Phi(Z)\|
  \le \|\Phi\|_F F(d(x,y)).
\end{equation}
\end{lemma}

\begin{proof}
This is the defining inequality obtained by multiplying
\cref{eq:F-norm} by $F(d(x,y))$.
\end{proof}

\begin{lemma}[Uniform tail]
Define

\begin{equation}
  T_F(R)
  =\sup_{x\in\Gamma}
  \sum_{\substack{y\in\Gamma\\d(x,y)\ge R}}F(d(x,y)).
\end{equation}

Then $T_F(R)\to0$ as $R\to\infty$ whenever the uniform summability in
\cref{eq:F-sum} is accompanied by the usual uniform tail property.  On a
uniformly locally finite lattice and for the standard polynomial or weighted
polynomial examples, this property holds.
\end{lemma}

\begin{proof}
For the standard lattice examples, compare the sum with the corresponding
radial series.  More generally, the tail property is recorded separately
because boundedness of a family of summable functions does not force uniform
integrability without a geometric hypothesis.  All later uniform convergence
statements assume $T_F(R)\to0$.
\end{proof}

\begin{remark}
This explicit tail assumption is useful in a family theorem.  It prevents a
silent exchange of a supremum over base points with a limit in the radius.
For $\mathbb Z^\nu$ and the $F$-functions used in applications, no extra work
is needed.
\end{remark}

\section{Parameterized interaction families}

\subsection{Uniform locality and local continuity}

Let $S$ be a compact Hausdorff space.  The intended applications include
profinite sets, compact coupling spaces, and compact disorder hulls.

\begin{definition}[Uniformly F-local family]
A family $\Phi=(\Phi_s)_{s\in S}$ is uniformly $F$-local if

\begin{equation}
  M_F(\Phi):=\sup_{s\in S}\|\Phi_s\|_F<\infty.
  \label{eq:family-M}
\end{equation}
\end{definition}

This is a property of the whole family.  It is not the statement that each
$\|\Phi_s\|_F$ is finite.

\begin{definition}[Local F-continuity]
A uniformly $F$-local family is locally $F$-continuous if, for every finite
$\Lambda\Subset\Gamma$, the map

\begin{equation}
  S\longrightarrow\mathcal B_F,
  \qquad
  s\longmapsto\Phi_s|_\Lambda
\end{equation}

is continuous, where $\Phi_s|_\Lambda$ retains only terms supported inside
$\Lambda$.  Equivalently, for every net $s_i\to s$ and finite $\Lambda$,

\begin{equation}
  \| (\Phi_{s_i}-\Phi_s)|_\Lambda\|_F\longrightarrow0.
  \label{eq:local-F-cont}
\end{equation}
\end{definition}

Since a finite $\Lambda$ has only finitely many subsets, local $F$-continuity
is equivalent to norm continuity of every coefficient map
$s\mapsto\Phi_s(Z)$ for finite $Z\Subset\Gamma$.  This coordinate
description is useful, but the restricted $F$-norm is better suited to the
estimates below.

Terms crossing the boundary of $\Lambda$ are not ignored analytically.  Their
effect is controlled uniformly by the $F$-tail after $\Lambda$ is enlarged.
This is why local continuity must be paired with
\cref{eq:family-M}.

\begin{proposition}[A sufficient global condition]
If $S$ is compact and $s\mapsto\Phi_s$ is continuous as a map into the
Banach space $\mathcal B_F$, then the family is uniformly $F$-local and
locally $F$-continuous.
\end{proposition}

\begin{proof}
The norm is continuous, so its image on the compact set $S$ is bounded.
Restriction to a finite volume is a continuous contraction in the relevant
interaction norm.
\end{proof}

\begin{remark}
The converse is false and should not be desired.  In a product disorder
space, two configurations that agree near the origin can disagree at every
distant scale.  A covariant interaction may therefore be close on every
fixed local region without being close in a norm that takes a supremum over
all lattice sites.
\end{remark}

\subsection{Finite-volume parameter continuity}

\begin{lemma}
If $\Phi$ is locally $F$-continuous, then for fixed finite $\Lambda$ the map

\begin{equation}
  (s,t,A)\longmapsto\tau_t^{\Lambda,\Phi_s}(A)
\end{equation}

is norm-continuous on $S\times\mathbb R\times\Alg_\Lambda$.
\end{lemma}

\begin{proof}
The finite sum $H_\Lambda^{\Phi_s}$ is norm-continuous in $s$.  The
exponential map is norm-continuous on finite-dimensional matrix algebras, and
conjugation is jointly continuous.
\end{proof}

This elementary lemma is not enough for the thermodynamic limit.  The number
of interaction terms grows with the volume.  Uniform locality is what makes
the finite statement stable as $\Lambda$ tends to $\Gamma$.

\section{Uniform Lieb-Robinson estimates}

\subsection{The finite-volume bound}

For finite $X\subseteq\Gamma$, let $\partial_{\Phi}X$ be the set of sites in
$X$ belonging to a nonzero interaction term, meaning a set $Z$ with
$\Phi(Z)\ne0$, that also meets $\Gamma\setminus X$.
For disjoint finite $X,Y$, define

\begin{equation}
  D_{F,\Phi}(X,Y)
  =\min\left\{
  \sum_{x\in X}\sum_{y\in\partial_\Phi Y}F(d(x,y)),
  \sum_{x\in\partial_\Phi X}\sum_{y\in Y}F(d(x,y))
  \right\}.
\end{equation}

We also use the coarser parameter-independent quantity

\begin{equation}
  \overline D_F(X,Y)
  =\sum_{x\in X}\sum_{y\in Y}F(d(x,y)).
  \label{eq:Dbar}
\end{equation}

Clearly $D_{F,\Phi}(X,Y)\le\overline D_F(X,Y)$.

\begin{theorem}[Lieb-Robinson bound]
\label{thm:LR-single}
Let $\Phi\in\mathcal B_F$.  Let $X,Y\Subset\Gamma$ be disjoint,
$\Lambda\supseteq X\cup Y$, $A\in\Alg_X$, and $B\in\Alg_Y$.  Then

\begin{equation}
\begin{split}
 \|[\tau_t^{\Lambda,\Phi}(A),B]\|
 \le {}&
 \frac{2\|A\|\|B\|}{C_F}
 \left(e^{2C_F\|\Phi\|_F|t|}-1\right)
 D_{F,\Phi}(X,Y).
 \label{eq:LR-single}
\end{split}
\end{equation}

In particular, the same estimate holds with $D_{F,\Phi}$ replaced by
$\overline D_F$.
\end{theorem}

\begin{proof}[Proof sketch with the analytic dependence displayed]
For $B\in\Alg_Y$, consider the commutator map
$A\mapsto[\tau_t^{\Lambda,\Phi}(A),B]$.  Differentiate after removing the
internal dynamics on $X$.  Only interaction terms crossing the current
support boundary enter the resulting integral inequality.  Iterating that
inequality produces paths of overlapping interaction supports.  The
$n$th-order term is bounded by

\begin{equation}
  2\|A\|\|B\|
  \frac{(2|t|\|\Phi\|_F)^n}{n!}
  C_F^{n-1}D_{F,\Phi}(X,Y).
\end{equation}

Summing over $n\ge1$ gives \cref{eq:LR-single}.  The convolution estimate
\cref{eq:F-conv} is exactly what turns the sum over intermediate lattice
sites into $C_F^{n-1}F(d(x,y))$.  A complete proof, including time-dependent
interactions and unbounded on-site terms, is Theorem 3.1 of
\cite{NSY2019}.
\end{proof}

\begin{theorem}[Uniform family Lieb-Robinson estimate]
\label{thm:family-LR}
Let $(\Phi_s)_{s\in S}$ be uniformly $F$-local with
$M=M_F(\Phi)$.  Under the hypotheses of \cref{thm:LR-single},

\begin{equation}
 \|[\tau_t^{\Lambda,\Phi_s}(A),B]\|
 \le
 \frac{2\|A\|\|B\|}{C_F}
 \left(e^{2C_FM|t|}-1\right)
 \overline D_F(X,Y)
 \label{eq:uniform-LR}
\end{equation}

for every $s\in S$ and every finite $\Lambda\supseteq X\cup Y$.
\end{theorem}

\begin{proof}
Apply \cref{thm:LR-single} to each $\Phi_s$, use
$\|\Phi_s\|_F\le M$, and replace the interaction-dependent boundary factor
by \cref{eq:Dbar}.  No compactness or topology on $S$ is used.
\end{proof}

\begin{remark}[Status]
The single-interaction estimate is \EST{}.  The family statement is a direct
corollary, also \EST{}.  The uniformity comes from the common norm bound, not
from the word ``family'' and not from profiniteness.
\end{remark}

\subsection{Exponential light-cone form}

Suppose $F_a(r)=e^{-ar}F(r)$ and the family is uniformly $F_a$-local.  Define

\begin{equation}
  v_a=\frac{2C_{F_a}M_{F_a}(\Phi)}{a}.
\end{equation}

\begin{corollary}
For disjoint $X,Y$,

\begin{equation}
\begin{split}
 \|[\tau_t^{\Lambda,\Phi_s}(A),B]\|
 \le{}&
 2\|A\|\|B\|\|F\|C_{F_a}^{-1}
 \min\{|X|,|Y|\} \\
 &\times
 \exp\bigl(a(v_a|t|-d(X,Y))\bigr),
 \label{eq:light-cone}
\end{split}
\end{equation}

uniformly in $s$ and $\Lambda$.
\end{corollary}

\begin{proof}
Subadditivity of distance extracts $e^{-ad(X,Y)}$ from the spatial sum.
Uniform summability bounds the remaining sum by
$\|F\|\min\{|X|,|Y|\}$.  Rewriting the time exponential with the definition
of $v_a$ gives \cref{eq:light-cone}.
\end{proof}

The number $v_a$ is a velocity bound, not necessarily the optimal physical
velocity.  The theorem controls a light-cone envelope and makes no assertion
that a signal saturates it.

\section{Thermodynamic dynamics}

\subsection{Existence and uniform convergence}

Let $(\Lambda_n)_{n\ge1}$ be increasing and exhaustive.  For a fixed local
observable $A\in\Alg_X$, choose $n_0$ with $X\subseteq\Lambda_{n_0}$.

\begin{theorem}[Uniform thermodynamic dynamics]
\label{thm:thermo}
Let $(\Phi_s)_{s\in S}$ be uniformly $F$-local, and assume the uniform
$F$-tail tends to zero.  For every $s\in S$, $A\in\Alg^{\mathrm{loc}}_\Gamma$,
and $t\in\mathbb R$, the limit

\begin{equation}
  \tau_t^s(A)=\lim_{n\to\infty}
  \tau_t^{\Lambda_n,\Phi_s}(A)
  \label{eq:thermo-limit}
\end{equation}

exists in norm.  It is independent of the exhaustion.  Convergence is
uniform in $s$ and in $t$ on compact intervals.  Each
$(\tau_t^s)_{t\in\mathbb R}$ extends uniquely to a strongly continuous group
of $*$-automorphisms of $\Alg_\Gamma$.
\end{theorem}

\begin{proof}
For $n\ge m\ge n_0$, the standard comparison estimate for two finite-volume
dynamics gives

\begin{equation}
\begin{split}
 &\|\tau_t^{\Lambda_n,\Phi_s}(A)
       -\tau_t^{\Lambda_m,\Phi_s}(A)\| \\
 &\qquad\le
 \frac{2\|A\|}{C_F}
 e^{2C_FM|t|}C_FM|t|
 \sum_{x\in X}
 \sum_{y\in\Lambda_n\setminus\Lambda_m}F(d(x,y)).
 \label{eq:Cauchy-tail}
\end{split}
\end{equation}

The right side is independent of $s$.  It tends to zero as $m,n\to\infty$,
uniformly for $t$ in a compact interval.  Thus the finite-volume dynamics
form a uniform Cauchy net on local observables.  The group, involution, norm,
and multiplicative properties pass to the limit.  Exhaustion independence
follows by interlacing two exhaustions.  Strong continuity holds first on the
dense local algebra using a uniform small-time estimate and then on all of
$\Alg_\Gamma$ by isometry.  This is the family-uniform version of Theorem
3.5 in \cite{NSY2019}.
\end{proof}

\begin{corollary}[Infinite-volume Lieb-Robinson bound]
The estimate \cref{eq:uniform-LR} holds with
$\tau_t^{\Lambda,\Phi_s}$ replaced by $\tau_t^s$.
\end{corollary}

\begin{proof}
Take the norm limit in the finite-volume commutator.  The bound is independent
of the volume.
\end{proof}

\subsection{Continuity in the family parameter}

\begin{theorem}[Point-norm parameter continuity]
\label{thm:param-cont}
Assume the hypotheses of \cref{thm:thermo} and local $F$-continuity.  For
every local $A$ and fixed $t$,

\begin{equation}
  S\longrightarrow\Alg_\Gamma,
  \qquad
  s\longmapsto\tau_t^s(A)
\end{equation}

is norm-continuous.  Continuity is uniform in $t$ on compact intervals.
\end{theorem}

\begin{proof}
Fix $s\in S$, a net $s_i\to s$, a local observable $A\in\Alg_X$, a compact
time interval $[-T,T]$, and $\epsilon>0$.  By
\cref{eq:Cauchy-tail}, choose a finite $\Lambda$ containing $X$ such that

\begin{equation}
  \sup_{r\in S,\,|t|\le T}
  \|\tau_t^r(A)-\tau_t^{\Lambda,\Phi_r}(A)\|
  <\epsilon/3.
\end{equation}

For this fixed $\Lambda$, local $F$-continuity implies
$H_\Lambda^{\Phi_{s_i}}\to H_\Lambda^{\Phi_s}$ in norm.  Finite-volume
parameter continuity then gives, eventually and uniformly for $|t|\le T$,

\begin{equation}
  \|\tau_t^{\Lambda,\Phi_{s_i}}(A)
       -\tau_t^{\Lambda,\Phi_s}(A)\|<\epsilon/3.
\end{equation}

The triangle inequality proves the claim.  This is also a specialization of
the local $F$-convergence theorem in \cite{NSY2019}.
\end{proof}

\begin{remark}
Global $\mathcal B_F$-continuity gives a direct quantitative estimate for the
difference of two dynamics.  The local argument above is weaker and better
adapted to disorder.  It works because a local observable sees a large but
finite region up to a uniformly controlled tail.
\end{remark}

\section{Profinite parameters and condensed organization}

\subsection{Why profinite spaces occur}

Let $D$ be finite.  The product

\begin{equation}
  \Omega=D^{\mathbb Z^d}
\end{equation}

is compact, Hausdorff, and totally disconnected, hence profinite.  Its clopen
cylinder sets record finite patches of a disorder configuration.  A local
rule assigning interaction terms from finite patches is continuous in the
product topology.  Uniform boundedness of the allowed coupling values can
give a common $F$-norm estimate, but this is an additional calculation.

Profinite sets also arise as inverse limits of finite-resolution parameter
spaces.  The useful fact is that continuous maps out of a profinite set can
be approximated, in suitable uniform targets, by locally constant maps.  For
interaction families, however, approximation must respect the uniform
$F$-bound.  Finite-resolution language does not excuse a missing tail
estimate.

\subsection{A functor of controlled interactions}

Fix $(\Gamma,d)$, the on-site algebras, and an $F$-function.  For $M>0$ and a
profinite set $S$, define

\begin{equation}
  \Ham_F^M(S)
  =\left\{
  (\Phi_s)_{s\in S}:
  \begin{array}{l}
  \|\Phi_s\|_F\le M\text{ for every }s,\\
  s\mapsto\Phi_s\text{ is locally }F\text{-continuous}
  \end{array}
  \right\}.
  \label{eq:Ham-functor}
\end{equation}

For a continuous map $f:T\to S$, define

\begin{equation}
  f^*:\Ham_F^M(S)\longrightarrow\Ham_F^M(T),
  \qquad
  (f^*\Phi)_t=\Phi_{f(t)}.
\end{equation}

\begin{proposition}[Functoriality]
The assignment $S\mapsto\Ham_F^M(S)$ is a contravariant functor on profinite
sets.  The thermodynamic dynamics are natural under pullback:

\begin{equation}
  \tau_t^{(f^*\Phi)_u}(A)=\tau_t^{\Phi_{f(u)}}(A).
\end{equation}
\end{proposition}

\begin{proof}
The norm bound is preserved by precomposition.  Local continuity is preserved
because a composite of continuous maps is continuous.  The dynamics identity
holds in every finite volume and therefore in the thermodynamic limit.
Identity maps and composites act as required.
\end{proof}

\begin{theorem}[Finite-cover descent]
\label{thm:descent}
On a fixed small site of profinite sets with finite jointly surjective covers,
$\Ham_F^M$ is a sheaf of sets.
\end{theorem}

\begin{proof}
Let $(S_i\to S)_{i=1}^k$ be a finite jointly surjective family and suppose
$\Phi_i\in\Ham_F^M(S_i)$ agree on every fiber product
$S_i\times_S S_j$.  For $s\in S$, choose a lift $s_i\in S_i$ and define
$\Phi_s=(\Phi_i)_{s_i}$.  Compatibility makes this independent of the lift.

The disjoint union $\coprod_i S_i\to S$ is a continuous surjection from a
compact space to a Hausdorff space, hence a quotient map.  Each finite-volume
coefficient map descends to a continuous map on $S$.  The common bound
$\|\Phi_s\|_F\le M$ is inherited fiberwise.  Thus the glued family belongs to
$\Ham_F^M(S)$.  Uniqueness follows from joint surjectivity.
\end{proof}

\begin{corollary}[Variable bounds]
Let

\begin{equation}
  \Ham_F^{\mathrm{ub}}(S)=\bigcup_{M>0}\Ham_F^M(S).
\end{equation}

Then $\Ham_F^{\mathrm{ub}}$ also satisfies finite-cover descent.
\end{corollary}

\begin{proof}
A finite cover supplies finitely many bounds $M_i$.  Their maximum is a global
bound for the glued family.
\end{proof}

\begin{remark}[What was proved]
The sheaf statement is \PROP{} as an organization of a familiar analytic
class.  Its proof is elementary topology.  The propagation theorem used to
define the controlled class remains the established analytic input from
Lieb-Robinson theory.
\end{remark}

\subsection{From a sheaf of sets to a moduli stack}

Interactions may have symmetries, gauge identifications, and higher
equivalences.  Replacing \cref{eq:Ham-functor} by a groupoid-valued or
anima-valued object is natural.  One might form an action groupoid for a
chosen group of quasi-local automorphisms, then impose descent after specifying
the topology on morphisms.

That construction is not carried out here.  In particular, the following
issues remain \OPEN{}:

\begin{itemize}[leftmargin=2.3em]
\item which automorphisms are morphisms before imposing a spectral gap;
\item how stabilization by product-state ancillas is represented;
\item whether a chosen localization preserves the analytic subfunctor;
\item how higher homotopies interact with a uniform locality witness;
\item which universe and hypercompletion conventions define the condensed
anima.
\end{itemize}

Our result supplies the object space and its dynamics.  It is not yet the full
condensed moduli stack.

\section{Examples}

\subsection{Uniformly bounded finite-range couplings}

Let $\Gamma=\mathbb Z^d$.  Fix an interaction range $R$ and a finite list of
shapes $Z_1,\ldots,Z_N$.  Suppose

\begin{equation}
  \Phi_s(x+Z_j)=J_j(s,x)K_{j,x},
\end{equation}

where $\|K_{j,x}\|\le1$ and

\begin{equation}
  \sup_{s,x,j}|J_j(s,x)|\le J_*.
\end{equation}

Only finitely many translates containing a given pair $x,y$ contribute, and
none contribute when $d(x,y)$ exceeds the maximal diameter of a shape.
Consequently,

\begin{equation}
  \sup_s\|\Phi_s\|_{F_a}
  \le C(d,R,N,F_a)J_*.
\end{equation}

If each $J_j(s,x)$ depends continuously on a finite cylinder of $s\in
D^{\mathbb Z^d}$, the family is locally $F$-continuous.  This gives a direct
class of profinite points of \cref{eq:Ham-functor}.

\subsection{Covariant disorder}

Let $\Omega=D^{\mathbb Z^d}$ with shift $T_x$.  A covariant local rule has the
form

\begin{equation}
  \Phi_\omega(x+Z)=U_x\Phi_{T_{-x}\omega}(Z)U_x^*,
\end{equation}

where $U_x$ translates observables.  A uniform bound on the finite set of
local coupling values gives uniform locality.  Product-topology continuity
gives local continuity.  The global map
$\omega\mapsto\Phi_\omega$ usually fails to be continuous in
$\mathcal B_F$, since the norm takes a supremum over all translations.

The family theorem is designed for precisely this situation.  It asks for the
uniform bound directly and uses local continuity only where continuity is
needed.

\subsection{Time-dependent interactions}

For completeness, let $u\mapsto\Phi_s(u)$ be strongly continuous in physical
time and locally bounded in $F$-norm.  On a compact time interval $I$, assume

\begin{equation}
  \sup_{s\in S}\int_I\|\Phi_s(u)\|_F\,du<\infty.
\end{equation}

The same proof yields a two-parameter cocycle $\tau^s_{t,r}$.  In the
time-independent exponent, replace $M|t-r|$ by

\begin{equation}
  \int_{\min(t,r)}^{\max(t,r)}\|\Phi_s(u)\|_F\,du,
\end{equation}

Thus $2C_FM|t-r|$ in the exponential becomes
$2C_F\int_{\min(t,r)}^{\max(t,r)}\|\Phi_s(u)\|_F\,du$.

This extension is \EST{} in \cite{NSY2019}.  We use time-independent
notation elsewhere to keep the parameter $s$ distinct from physical time.

\section{Counterexamples and sharp distinctions}

\subsection{Pointwise finite norms do not give a common norm}

\begin{counterexample}
Let $S=\mathbb N\cup\{\infty\}$ be the one-point compactification and let
$\Psi$ be a nonzero finite-range interaction.  Put

\begin{equation}
  \Phi_n=n\Psi,
  \qquad
  \Phi_\infty=0.
\end{equation}

Every fiber has finite $F$-norm, but
$\sup_{s\in S}\|\Phi_s\|_F=\infty$.
\end{counterexample}

This family is not locally continuous at infinity, so it is intentionally
simple.  It shows that a pointwise finiteness statement has the wrong
quantifiers for uniform propagation.

\subsection{Coefficientwise continuity needs a uniform tail condition}

One can arrange interaction mass at increasing distances so that every fixed
coefficient stabilizes while a global norm diverges.  For example, choose
pair terms on $\mathbb Z$ supported at $\{0,n\}$ with amplitudes tuned so that
each interaction has finite norm but the norms grow with $n$.  On the
one-point compactification, coefficients on each fixed finite support converge
to zero.  No common Lieb-Robinson velocity follows.

This is the reason our local continuity definition is not presented as a
substitute for \cref{eq:family-M}.  It is a second condition with a different
job.

\subsection{The moving weak-defect counterexample}

The next example is central because it satisfies uniform locality and local
continuity while its gaps collapse.

\begin{counterexample}[Moving weak defect]
\label{cex:moving-defect}
Let $\Gamma=\mathbb N$ with one qubit at each site.  Write

\begin{equation}
  p_x=|1\rangle\langle1|_x.
\end{equation}

For $n\ge1$, define the on-site interaction

\begin{equation}
  \Phi_n(\{x\})=
  \begin{cases}
  n^{-1}p_n,&x=n,\\
  p_x,&x\ne n,
  \end{cases}
  \qquad
  \Phi_n(Z)=0\quad\text{if }|Z|\ne1.
  \label{eq:defect-n}
\end{equation}

Define the limiting interaction by

\begin{equation}
  \Phi_\infty(\{x\})=p_x.
  \label{eq:defect-infty}
\end{equation}

Let $S=\mathbb N\cup\{\infty\}$.  Then:

\begin{enumerate}[leftmargin=2.3em]
\item the family is finite-range and uniformly $F$-local;
\item $\Phi_n\to\Phi_\infty$ locally in $F$-norm;
\item all fibers have the same unique product ground state;
\item the GNS gaps are $\gamma_n=1/n$ and $\gamma_\infty=1$;
\item therefore every fiber is gapped but
$\inf_{s\in S}\gamma_s=0$.
\end{enumerate}
\end{counterexample}

\begin{proof}
Only on-site terms occur, so the interaction range is zero.  With the usual
normalization $F(0)>0$,

\begin{equation}
  \|\Phi_n\|_F
  =F(0)^{-1}\sup_x\|\Phi_n(\{x\})\|
  =F(0)^{-1}.
\end{equation}

Thus the family has a common norm bound.  For a fixed finite
$\Lambda\subset\mathbb N$, the defect site $n$ eventually lies outside
$\Lambda$, so $\Phi_n|_\Lambda=\Phi_\infty|_\Lambda$.  This proves local
$F$-convergence.

The vector with every qubit in $|0\rangle$ is the unique zero-energy product
ground state.  In the $n$th fiber, flipping the qubit at site $n$ costs
$1/n$, while every other single flip costs $1$.  All multi-flip energies are
sums of positive on-site costs.  Hence the GNS gap is $1/n$.  In the limiting
fiber every single flip costs $1$, so its gap is $1$.
More explicitly, the GNS Hamiltonian is diagonal on the dense basis of
finite-spin-flip vectors.  Its positive eigenvalues are finite sums of the
on-site costs.  Their infimum is therefore the least on-site cost, so no
infrared sequence produces an additional spectral value below the stated
gap.
\end{proof}

\begin{remark}
The limiting interaction in \cref{eq:defect-infty} does not omit the defect
term.  Instead, the weak term moves away and the local limit restores unit
strength at every fixed site.  This choice ensures that the limiting fiber is
also gapped.  It makes the counterexample stronger than one whose limiting
fiber is simply gapless.
\end{remark}

\begin{corollary}
Uniform Lieb-Robinson estimates, local parameter continuity, and pointwise
gappedness do not imply a parameter-uniform gap.
\end{corollary}

\begin{proof}
The first two properties and failure of the third follow from
\cref{cex:moving-defect}.
\end{proof}

\subsection{Compact norm-continuous matrices and changing degeneracy}

There is also a finite-dimensional obstruction.  Let

\begin{equation}
  S=\{0\}\cup\{1/n:n\ge2\},
  \qquad
  H_s=\operatorname{diag}(0,s,1).
\end{equation}

This is a norm-continuous family on a profinite compact set.  At $s=0$, the
literal ground space has dimension two and the gap above it is $1$.  At
$s=1/n$, the ground state is unique and the literal ground-state gap is
$1/n$.  Every fiber is gapped, but there is no common lower bound.

If the two lowest levels are instead treated as one isolated ground band, the
gap from that band to the level at $1$ is at least $1/2$.  Thus a later
definition of a uniformly gapped substack must specify whether it tracks the
literal ground eigenspace or an isolated low-energy band of constant rank.

\section{Finite approximation and executable diagnostics}

The companion Haskell program performs finite calculations.  It is useful
for checking constants and quantifiers, but it is not a proof assistant and
does not establish a thermodynamic theorem.

\subsection{Finite F-function diagnostics}

On a finite line $\{0,\ldots,L-1\}$, the program computes

\begin{equation}
  \|F\|_L
  =\max_x\sum_yF(|x-y|)
\end{equation}

and

\begin{equation}
  C_{F,L}
  =\max_{x,y}
  \frac{\sum_zF(|x-z|)F(|z-y|)}{F(|x-y|)}.
\end{equation}

For increasing $L$, these values provide finite diagnostics for candidate
$F$-functions.  Bounded values in a numerical sample do not prove boundedness
for the infinite lattice.  An analytic estimate remains necessary.

\subsection{Lieb-Robinson envelope}

Given finite estimates $C_{F,L}$, $M$, supports $X,Y$, and time $t$, the code
evaluates

\begin{equation}
  \mathcal E_L(t;X,Y)
  =\frac{2}{C_{F,L}}
  \left(e^{2C_{F,L}M|t|}-1\right)
  \sum_{x\in X,y\in Y}F(|x-y|).
\end{equation}

Norms of $A$ and $B$ are set separately.  This calculation is useful for
catching missing factors and for seeing how a chosen decay profile affects
the envelope.

\subsection{Pointwise versus uniform predicates}

For a finite list of positive gaps, both ``every gap is positive'' and
``there is a positive minimum'' hold.  The distinction appears in a family
of growing test sets.  For the defect sequence,

\begin{equation}
  \min_{1\le n\le N}\gamma_n=\frac1N.
\end{equation}

The program prints this value for increasing $N$.  The analytic formula, not
the finite run, proves that the infimum over all $n$ is zero.

\section{Interfaces required by later papers}

\subsection{Input to the gap paper}

The thermodynamic gap is defined using the generator of the infinite-volume
dynamics in a chosen ground-state representation.  Paper 3 therefore needs
the following outputs from the present paper:

\begin{enumerate}[leftmargin=2.3em]
\item the quasi-local algebra $\Alg_\Gamma$;
\item the class of admissible interactions $\mathcal B_F$;
\item a common locality witness $M_F(\Phi)$ for a family;
\item the thermodynamic automorphism group $\tau_t^s$;
\item point-norm parameter continuity on local observables;
\item the distinction between local and global interaction topologies.
\end{enumerate}

Paper 3 must add a ground-state convention and a gap witness.  No gap
predicate is part of \cref{thm:family-LR} or \cref{thm:thermo}.

\subsection{Input to spectral flow}

Quasi-adiabatic continuation assumes a differentiable interaction path, a
uniform isolated spectral band, and stronger decay control on the interaction
and its derivative.  Under those hypotheses, Bachmann, Michalakis,
Nachtergaele, and Sims construct a quasi-local spectral-flow automorphism
\cite{BMNS2012}.

The present paper supplies the propagation and thermodynamic-limit pattern
used by that construction.  It does not verify the spectral gap along a path.
Accordingly, the implication

\begin{equation}
  \text{uniformly local path}
  \Longrightarrow
  \text{spectral-flow equivalence}
\end{equation}

is \OBS{} without an independent uniform-gap assumption.

\subsection{Input to observable algebras and K-theory}

Paper 2 begins with the same quasi-local algebra and may form
$C(S,\Alg_\Gamma)$ for a compact parameter space $S$.  The existence of this
C*-algebra of continuous functions is independent of the existence of a
Hamiltonian dynamics.  Conversely, the $F$-norm contains interaction data not
recoverable from the abstract C*-algebra alone.

Thus the arrows

\begin{equation}
  \Phi_s\longmapsto\Alg_\Gamma
  \longmapsto C(S,\Alg_\Gamma)
\end{equation}

forget information.  Any later K-theory invariant should be described as an
invariant of the chosen observable-algebra construction, not as a complete
encoding of the interaction.

\section{Status ledger and limitations}

\begin{longtable}{>{\raggedright\arraybackslash}p{0.26\textwidth}
                  >{\raggedright\arraybackslash}p{0.18\textwidth}
                  >{\raggedright\arraybackslash}p{0.47\textwidth}}
\toprule
Statement & Status & Reason \\
\midrule
\endfirsthead
\toprule
Statement & Status & Reason \\
\midrule
\endhead
Finite-volume Lieb-Robinson estimate under finite $F$-norm
& \EST{}
& Theorem 3.1 of \cite{NSY2019}, with earlier foundations in
\cite{NachtergaeleSims2010}. \\
Uniform estimate for an $S$-family with a common $F$-norm bound
& \EST{}
& Direct substitution of the common bound into the cited estimate. \\
Uniform thermodynamic dynamics
& \EST{}
& Uniform version of the Cauchy-tail proof in \cite{NSY2019}. \\
Parameter continuity from local $F$-continuity and a common tail bound
& \EST{}
& Finite truncation plus uniform tails, compatible with the local
$F$-convergence theorem of \cite{NSY2019}. \\
The functor $\Ham_F^{\mathrm{ub}}$ on profinite sets
& \PROP{}
& Definition introduced here to organize the established analytic class. \\
Finite-cover descent for $\Ham_F^{\mathrm{ub}}$
& \EST{}
& Elementary gluing over quotient maps, once the analytic class is fixed. \\
Canonical condensed anima of all Hamiltonians and equivalences
& \OPEN{}
& Morphisms, higher coherences, topology, stabilization, and size
conventions remain to be specified. \\
Uniform locality from profiniteness alone
& \OBS{}
& Profinite topology does not bound interaction norms. \\
Uniform gap from uniform locality and pointwise gaps
& \OBS{}
& The moving weak-defect family in \cref{cex:moving-defect}. \\
\bottomrule
\end{longtable}

\subsection{Analytic limitations}

We work with bounded interaction terms and finite-dimensional on-site spaces.
Bosonic lattice models, continuum systems, and interactions with unbounded
terms require different domains, representations, or state-dependent bounds.
The conclusions should not be transferred to those settings by analogy.

The $F$-norm is sufficient, not necessary.  Some models admit sharper or
state-dependent propagation estimates even when the norm used here is
infinite.  Our functor therefore describes a robust controlled subspace of
interactions, not every physically local model.

The light-cone estimate is an upper bound.  It does not identify the exact
front velocity, diffusion law, or transport coefficient.  Disorder can
produce anomalous or zero-velocity bounds that are stronger than the common
estimate used here.

\subsection{Categorical limitations}

The sheaf in \cref{thm:descent} is set-valued.  It does not encode unitary
conjugacies, quasi-local automorphisms, defects, or higher families.  Calling
it a moduli stack without adding those morphisms would overstate the result.

The finite-cover sheaf property does not say that every invariant factors
through a finite quotient of a profinite parameter.  Locally constant
families may approximate continuous families in a chosen norm, but an
invariant must also be continuous in that norm before an approximation
argument applies.

The associated condensed object depends on the selected topology on the
interaction class.  The global $F$-norm topology is convenient for analysis
but too strong for many disorder families.  The local $F$-topology admits
those families but does not make the norm function locally bounded without a
separate uniformity condition.  This tradeoff is structural, not cosmetic.

\subsection{Physical limitations}

A Lieb-Robinson bound controls propagation of commutators.  It does not imply
a ground state, uniqueness, positivity of a gap, topological order, or
invertibility.  It also does not distinguish a trivial phase from a
topological one.

The moving defect example shows that even an exceptionally simple family of
commuting on-site Hamiltonians can violate uniform gappedness.  Any proposed
gapped substack must therefore carry the uniform gap as data or as a proved
property with the correct quantifiers.

\section{Discussion}

The condensed perspective is most useful here as bookkeeping with teeth.  It
forces one to ask what an $S$-family is, which topology controls variation,
and whether a statement is stable under pullback and finite-cover descent.
Those questions expose a hidden ambiguity in the slogan ``continuous family
of local Hamiltonians.''  Continuity and locality are not single predicates.

The global interaction Banach topology makes compactness powerful.  A
continuous map from a compact $S$ has a bounded image, so uniform propagation
follows immediately.  The same topology is poorly matched to covariant
disorder because it monitors every translate at once.  The local topology is
better matched to physical observables, but compactness no longer supplies a
global norm bound.  The correct definition records both local continuity and
uniform locality.

This two-part definition also behaves well under a finite profinite cover.
Local continuity glues because continuous maps glue across a compact quotient.
The uniform bound glues because a finite maximum exists.  That is the precise
sense in which condensed organization helps.  It preserves a carefully
defined analytic class under the covers used by the site.

The next step of the larger research program is harder.  To define a
uniformly gapped subfunctor, one must choose between finite-volume and GNS
gaps, state the boundary conditions, track a ground band, and demand one gap
constant for the whole parameter object.  The counterexamples here show that
none of this data can be suppressed.

There is also a useful lesson for formalization.  A proof assistant interface
should not encode ``local interaction'' as an unstructured label.  It should
carry an $F$-function, proofs of summability and convolution, an interaction,
and a numerical norm witness.  A family should carry a common witness rather
than a function returning a separate witness for every point.  That design
makes the quantifier distinction visible in the type.

\section{Conclusion}

We have established a controlled first arrow

\begin{equation}
  \text{local quantum interactions}
  \longrightarrow
  \text{parameterized thermodynamic dynamics}.
\end{equation}

The analytic theorem is simple to state: a common interaction $F$-norm bound
gives common Lieb-Robinson constants and uniform thermodynamic convergence.
Local $F$-continuity then gives continuity of the dynamics on local
observables.  These results are inherited from established Lieb-Robinson
theory with the family quantifiers made explicit.

We also proposed a profinite functor of controlled interaction families and
proved finite-cover descent.  This provides an honest object-level input for
a future condensed moduli stack.  Its categorical form does not produce the
analytic estimates.  It records the class for which those estimates have
already been proved.

Finally, the moving weak defect shows why the next arrow cannot be automatic.
Uniform locality and pointwise gaps coexist with a collapsing family gap.
The uniformly gapped substack must therefore impose new analytic data.  This
boundary between organization and analysis is the main conclusion of the
paper.

\appendix

\section{Detailed comparison estimate}

This appendix records the finite-volume comparison used in
\cref{thm:thermo}.  Let $\Lambda_m\subseteq\Lambda_n$ and write

\begin{equation}
  H_{\Lambda_n}=H_{\Lambda_m}+V_{m,n},
\end{equation}

where $V_{m,n}$ contains terms not wholly supported in $\Lambda_m$.  Duhamel's
formula gives

\begin{equation}
\begin{split}
 &\tau_t^{\Lambda_n}(A)-\tau_t^{\Lambda_m}(A)\\
 &\quad=
 i\int_0^t
 \tau_r^{\Lambda_n}
 \left([V_{m,n},\tau_{t-r}^{\Lambda_m}(A)]\right)dr.
 \label{eq:duhamel}
\end{split}
\end{equation}

Taking norms and decomposing $V_{m,n}$ into interaction terms yields

\begin{equation}
 \|\tau_t^{\Lambda_n}(A)-\tau_t^{\Lambda_m}(A)\|
 \le
 \sum_{Z\in\mathcal S_{m,n}}
 \int_0^{|t|}
 \|[\Phi(Z),\tau_{|t|-r}^{\Lambda_m}(A)]\|dr,
\end{equation}

where $\mathcal S_{m,n}$ is the set of new terms meeting the complement of
$\Lambda_m$.  The Lieb-Robinson estimate bounds each commutator by a spatial
sum from $X$ to $Z$.  Summing over $Z$ and using the interaction norm gives a
tail from $X$ to $\Lambda_n\setminus\Lambda_m$.  Integrating the time
exponential gives \cref{eq:Cauchy-tail}, up to the harmless choice between
$e^{2C_FM|t|}-1$ and its upper bound
$2C_FM|t|e^{2C_FM|t|}$.

Every constant in this argument is controlled by $F$, $C_F$, $M$, $|X|$,
and the time interval.  This is the entire reason the thermodynamic limit is
uniform in the family parameter.

\section{A finite-cover gluing lemma}

\begin{lemma}
Let $q:K\to S$ be a continuous surjection from a compact space to a Hausdorff
space.  Let $Y$ be any topological space.  A map $f:S\to Y$ is continuous if
and only if $f\circ q$ is continuous.
\end{lemma}

\begin{proof}
The map $q$ is closed because images of compact sets are compact and compact
subsets of a Hausdorff space are closed.  A continuous closed surjection is a
quotient map.  The assertion is the defining property of the quotient
topology.
\end{proof}

\begin{remark}
In \cref{thm:descent}, take $K=\coprod_iS_i$.  The target $Y$ is the
finite-volume interaction space for each fixed $\Lambda$.  The lemma proves
continuity of the glued local restrictions.  The analytic norm bound is
checked separately.
\end{remark}

\section{Logical form of uniform conditions}

The following formulas summarize the quantifiers that the paper keeps
separate.

Pointwise locality is

\begin{equation}
  \forall s\in S\;\exists M_s<\infty:
  \|\Phi_s\|_F\le M_s.
\end{equation}

Uniform locality is

\begin{equation}
  \exists M<\infty\;\forall s\in S:
  \|\Phi_s\|_F\le M.
\end{equation}

Pointwise gappedness is

\begin{equation}
  \forall s\in S\;\exists\gamma_s>0:
  \operatorname{gap}(H_s)\ge\gamma_s.
\end{equation}

Uniform gappedness is

\begin{equation}
  \exists\gamma>0\;\forall s\in S:
  \operatorname{gap}(H_s)\ge\gamma.
\end{equation}

Neither exchange of quantifiers is a matter of notation.  Compactness permits
the exchange only when the relevant bound or gap function has suitable
continuity properties.  The ground-state gap can fail to be lower
semicontinuous when ground-state multiplicity changes.  The moving defect
shows failure even with a fixed unique product ground state if the interaction
topology is local.

\section{Notation and reproducibility}

\begin{longtable}{>{\raggedright\arraybackslash}p{0.24\textwidth}
                  >{\raggedright\arraybackslash}p{0.66\textwidth}}
\toprule
Symbol & Meaning \\
\midrule
\endfirsthead
\toprule
Symbol & Meaning \\
\midrule
\endhead
$\Gamma$ & Countable metric lattice or graph. \\
$\Alg_\Lambda$ & Matrix algebra of observables in finite volume $\Lambda$. \\
$\Alg_\Gamma$ & Norm-completed quasi-local C*-algebra. \\
$\Phi$ & Bounded self-adjoint interaction. \\
$F$ & Uniformly summable decay function with a convolution bound. \\
$C_F$ & Convolution constant of $F$. \\
$\|\Phi\|_F$ & Interaction norm in \cref{eq:F-norm}. \\
$S$ & Compact parameter space, often profinite. \\
$M_F(\Phi)$ & Uniform interaction-norm bound over $S$. \\
$\tau_t^{\Lambda,\Phi_s}$ & Finite-volume dynamics in fiber $s$. \\
$\tau_t^s$ & Thermodynamic dynamics in fiber $s$. \\
$\Ham_F^M$ & Proposed sheaf of locally continuous families with common bound $M$. \\
\bottomrule
\end{longtable}

The Haskell sources accompanying this article are in the repository directory

\begin{center}
\path{src/locality-condensed-families/}.
\end{center}

They compile with GHC using
\texttt{-Wall}, \texttt{-Wextra}, and \texttt{-Werror}.  All exported functions have explicit type
signatures.  The computations are deterministic and require no external
packages beyond \texttt{base}.

\begin{thebibliography}{99}

\bibitem{LiebRobinson1972}
E. H. Lieb and D. W. Robinson,
The finite group velocity of quantum spin systems,
\emph{Communications in Mathematical Physics} \textbf{28} (1972), 251 to 257.
\href{https://doi.org/10.1007/BF01645779}{doi:10.1007/BF01645779}.

\bibitem{NachtergaeleSims2006}
B. Nachtergaele and R. Sims,
Lieb-Robinson bounds and the exponential clustering theorem,
\emph{Communications in Mathematical Physics} \textbf{265} (2006), 119 to 130.
\href{https://arxiv.org/abs/math-ph/0506030}{arXiv:math-ph/0506030}.

\bibitem{NachtergaeleOgataSims2006}
B. Nachtergaele, Y. Ogata, and R. Sims,
Propagation of correlations in quantum lattice systems,
\emph{Journal of Statistical Physics} \textbf{124} (2006), 1 to 13.
\href{https://arxiv.org/abs/math-ph/0603064}{arXiv:math-ph/0603064}.

\bibitem{NachtergaeleSims2010}
B. Nachtergaele and R. Sims,
Lieb-Robinson bounds in quantum many-body physics,
in \emph{Entropy and the Quantum}, Contemporary Mathematics \textbf{529}
(2010), 141 to 176.
\href{https://arxiv.org/abs/1004.2086}{arXiv:1004.2086}.

\bibitem{NVZ2011}
B. Nachtergaele, A. Vershynina, and V. A. Zagrebnov,
Lieb-Robinson bounds and existence of the thermodynamic limit for a class of
irreversible quantum dynamics,
in \emph{Entropy and the Quantum II}, Contemporary Mathematics \textbf{552}
(2011), 161 to 175.
\href{https://arxiv.org/abs/1103.1122}{arXiv:1103.1122}.

\bibitem{BMNS2012}
S. Bachmann, S. Michalakis, B. Nachtergaele, and R. Sims,
Automorphic equivalence within gapped phases of quantum lattice systems,
\emph{Communications in Mathematical Physics} \textbf{309} (2012), 835 to 871.
\href{https://arxiv.org/abs/1102.0842}{arXiv:1102.0842}.
\href{https://doi.org/10.1007/s00220-011-1380-0}{doi:10.1007/s00220-011-1380-0}.

\bibitem{NSY2019}
B. Nachtergaele, R. Sims, and A. Young,
Quasi-locality bounds for quantum lattice systems. Part I.
Lieb-Robinson bounds, quasi-local maps, and spectral flow automorphisms,
\emph{Journal of Mathematical Physics} \textbf{60} (2019), 061101.
\href{https://arxiv.org/abs/1810.02428}{arXiv:1810.02428}.
\href{https://doi.org/10.1063/1.5095769}{doi:10.1063/1.5095769}.

\bibitem{ClausenScholze}
D. Clausen and P. Scholze,
Condensed mathematics,
lecture notes, 2019.
\url{https://www.math.uni-bonn.de/people/scholze/Condensed.pdf}.

\end{thebibliography}

\end{document}
